% experiments.tex
\section{Experiments}
\label{sec:experiments}

We evaluate INSTRUMENTAL through four axes: (1)~the target sound and evaluation protocol, (2)~an ablation study over the synthesizer parameter space, (3)~a suite of hypothesis tests probing alternative synthesis modules and loss functions, and (4)~a comparison of optimizers.
All experiments use 44.1\,kHz audio on an Apple M4 Mini (10 CPU cores, 16\,GB unified memory) unless otherwise noted.

%% ─────────────────────────────────────────────────
\subsection{Target Sound and Evaluation Protocol}
\label{sec:target-sound}

The target is a lead synthesizer extracted from a commercial song.
We isolate 22 individual notes at three pitches---$\text{A3}=221$\,Hz, $\text{C\#4}=278$\,Hz, and $\text{D4}=295$\,Hz---spanning approximately 3.2 seconds at ${\sim}100$\,BPM, with each note lasting ${\sim}150$\,ms.
The notes are segmented using onset detection and pitch tracking via \texttt{librosa.pyin}.

For optimization, we select a representative subset of three notes (one per pitch class) as fitting targets.
The loss $\mathcal{L}$ averaged over these three notes serves as the primary metric throughout all experiments.
We use the \emph{matching loss}---a weighted combination of perceptually-scaled multi-resolution STFT, spectral centroid distance, and MFCC distance:
\begin{equation}
  \mathcal{L}_{\text{match}} = 1.0 \cdot \mathcal{L}_{\text{MelSTFT}} + 0.1 \cdot \mathcal{L}_{\text{centroid}} + 0.05 \cdot \mathcal{L}_{\text{MFCC}}
  \label{eq:matching-loss}
\end{equation}
where $\mathcal{L}_{\text{MelSTFT}}$ is a multi-resolution mel-scaled STFT loss with perceptual weighting (FFT sizes 1024, 2048, 8192), $\mathcal{L}_{\text{centroid}}$ is the normalized spectral centroid $L_1$ distance, and $\mathcal{L}_{\text{MFCC}}$ is the MSE over 13 MFCCs.
Full-sequence evaluation renders all 21 notes with the recovered parameters and compares against the original via RMS-normalized waveform alignment.

%% ─────────────────────────────────────────────────
\subsection{Ablation Study: Synthesizer Parameterization}
\label{sec:ablation}

We incrementally expand the synthesizer's parameter space to determine the minimal representation that captures the target timbre.
Each configuration is optimized with CMA-ES for 100K evaluations (population size~48, $\sigma_0 = 0.4$) using spectral initialization and 8-core CPU parallelism.

\begin{table}[ht]
\centering
\caption{Ablation study over synthesizer parameter count. Loss is the matching loss averaged over three target notes (lower is better).}
\label{tab:ablation}
\begin{tabular}{@{}clcc@{}}
\toprule
\textbf{Params} & \textbf{Modules added} & \textbf{Loss} & \textbf{$\Delta$} \\
\midrule
15 & Osc.\ mix, filter, ADSR, reverb & 4.60 & --- \\
18 & + unison voices/spread, noise floor & 2.34 & $-49.1\%$ \\
24 & + pulse width, filter ADSR, filter slope & 2.13 & $-9.0\%$ \\
28 & + 2-band parametric EQ & 2.09 & $-1.9\%$ \\
29 & + distortion, delay, vibrato (unconstrained) & \multicolumn{2}{c}{\emph{diverged}} \\
\bottomrule
\end{tabular}
\end{table}

The largest improvement comes from adding unison and noise floor (15$\to$18 params, $-49\%$), which allows the optimizer to model the characteristic ``thickness'' of the target.
Adding pulse width and a dedicated filter envelope (18$\to$24) yields a further 9\% reduction.
The 2-band parametric EQ (24$\to$28) provides a modest additional gain.

At 29 parameters, including unconstrained distortion drive, delay feedback, and vibrato depth, the search \emph{diverged}: CMA-ES exploited the extended parameter ranges, converging to physically implausible configurations with extreme delay feedback and distortion.
This motivated constraining the parameter space to 24 core synthesis parameters plus 4 post-EQ parameters (28 total), which we adopt as our final configuration.

%% ─────────────────────────────────────────────────
\subsection{Hypothesis Testing}
\label{sec:hypotheses}

We test eight hypotheses about alternative synthesis modules, loss functions, and optimization strategies.
Each test modifies one component while holding all others fixed.
All tests use CMA-ES with a budget of 10K evaluations (population size~20, $\sigma_0 = 0.15$), initialized from the best 24-parameter solution (v24, baseline loss $= 2.13$).
Table~\ref{tab:hypotheses} summarizes the results.

\begin{table}[ht]
\centering
\caption{Hypothesis tests. Each row modifies one component relative to the v24 baseline ($\mathcal{L} = 2.13$). Positive $\Delta$ means worse; negative means better. All tests use 10K CMA-ES evaluations.}
\label{tab:hypotheses}
\begin{tabular}{@{}llccc@{}}
\toprule
\textbf{ID} & \textbf{Hypothesis} & \textbf{Extra params} & \textbf{Loss} & \textbf{$\Delta$} \\
\midrule
A & Chebyshev waveshaping (5-coeff.) & +5 & 2.75 & $+29\%$ \\
B & 2-band parametric EQ & +4 & 1.96 & $-8\%$ \\
C & TFS / envelope loss & 0 & 2.49 & $+17\%$ \\
D & CDPAM perceptual loss & 0 & \multicolumn{2}{c}{\emph{skipped}\footnotemark} \\
E & SPSA fine-tuning (5K CMA + 5K SPSA) & 0 & 2.36 & $+11\%$ \\
F & CLAP embedding loss & 0 & 0.095\textsuperscript{*} & --- \\
G & Multi-start CMA-ES ($8 \times 1.25$K) & 0 & 2.75 & $+29\%$ \\
H & FM synthesis (carrier + modulator) & +2 & 2.49 & $+17\%$ \\
\bottomrule
\end{tabular}
\end{table}
\footnotetext{CDPAM skipped due to unresolved dependency conflict at experiment time.}

\noindent\textsuperscript{*}CLAP loss converged in its own metric space (cosine distance 0.095) but the resulting audio had perceptually different quality from the spectral-loss solution, as CLAP embeddings capture high-level semantic similarity rather than fine spectral detail.

\paragraph{Key findings.}
\begin{itemize}
  \item \textbf{Parametric EQ (B)} is the only modification that \emph{improves} on the baseline ($-8\%$), confirming that post-synthesis spectral shaping captures residual timbral detail the core oscillator-filter chain cannot model. This result motivated its inclusion in the final 28-parameter configuration (Section~\ref{sec:ablation}).
  \item \textbf{Chebyshev waveshaping (A)} and \textbf{FM synthesis (H)} increase loss by 29\% and 17\%, respectively. Both add harmonic complexity that the 10K-evaluation budget cannot adequately explore, suggesting these modules require either a larger search budget or a reduced base parameter set to be effective.
  \item \textbf{TFS/envelope loss (C)} adds an RMS-envelope trajectory term to the matching loss. The resulting parameters score 17\% worse under the standard loss, indicating that the envelope term distorts the spectral optimization landscape.
  \item \textbf{SPSA fine-tuning (E)} splits the budget into 5K CMA-ES + 5K SPSA gradient steps. The two-phase approach performs 11\% worse, likely because SPSA's noisy gradient estimates are too imprecise in the 24-dimensional space to improve on CMA-ES alone at this budget.
  \item \textbf{Multi-start CMA-ES (G)} distributes the same 10K budget across 8 independent restarts ($8 \times 1.25$K each). This performs 29\% worse, confirming that with a warm-start initialization the landscape is unimodal enough that concentrating the full budget on a single run is preferable.
  \item \textbf{CLAP loss (F)} converges rapidly but in a different quality regime---the CLAP embedding space captures ``sounds like a synthesizer'' rather than ``has the same spectral envelope,'' making it unsuitable as a drop-in replacement for spectral losses in parameter recovery tasks.
\end{itemize}

%% ─────────────────────────────────────────────────
\subsection{Optimizer Comparison}
\label{sec:optimizer-comparison}

We compare three optimization approaches on the 24-parameter synthesizer with matching loss.

\paragraph{CMA-ES.}
Our primary optimizer. Population size 48--80, spectral initialization, 100K evaluations.
CMA-ES converges to loss ${\approx}\,2.09$ (28 params) and ${\approx}\,2.13$ (24 params), consistently finding high-quality solutions across runs.

\paragraph{Adam (gradient descent).}
We use Adam with cosine-annealing learning rate ($\text{lr}_0 = 0.02$, $\eta_{\min} = 0.001$) and gradient noise injection for exploration, running for 500 steps with the differentiable synthesizer.
Adam plateaus at a higher loss than CMA-ES on the real-audio target: the multi-resolution STFT landscape contains many shallow local minima that trap gradient-based search.
On \emph{synthetic} targets (Level~0: known ground truth), Adam converges reliably to loss $< 0.01$ in under 500 steps, confirming that the synth and loss are fully differentiable---the difficulty lies in the real-audio matching landscape, not the optimization machinery.

\paragraph{RL--PPO.}
A PPO agent (2-layer MLP, 256 units, 8 parallel environments) trained for 1M timesteps learns to adjust 18 synth parameters over 50-step episodes.
The state space comprises current parameters (18-dim), target audio features (17-dim MFCC + centroid + rolloff + flatness + RMS), current audio features, and their difference (total: 69-dim).
Actions are continuous adjustments in $[-0.1, +0.1]$ per parameter.
The agent matches CMA-ES quality on synthetic targets but does not surpass it on real audio, and requires hours of training versus seconds for CMA-ES inference.
The RL approach demonstrates that sound matching \emph{can} be framed as sequential decision-making, but the current reward signal (loss improvement per step) provides insufficient shaping for the agent to learn fine spectral adjustments.

\paragraph{Summary.}
CMA-ES with warm-start initialization remains the best approach for single-target parameter recovery: it handles the non-convex landscape without requiring gradients, scales adequately to 24--28 parameters, and produces the lowest loss among all tested optimizers.

%% ─────────────────────────────────────────────────
\subsection{Compute Performance}
\label{sec:compute}

Throughput is critical for evolutionary search: CMA-ES evaluates thousands of candidates per generation.
We benchmark two implementation strategies.

\paragraph{Batched CPU synthesis.}
The GPU-style batched synthesizer (\texttt{SynthPatchGPU}) processes an entire CMA-ES population ($B = 80$ candidates) in a single batched forward pass, with all oscillator, filter, and envelope computations expressed as $(B, N)$ tensor operations.
On the M4 Mini (10 CPU cores), this achieves \textbf{553 evaluations per second}, rendering 80 candidate waveforms of ${\sim}150$\,ms each in one pass.

\paragraph{Multi-process CPU synthesis.}
An alternative 8-process \texttt{multiprocessing.Pool} implementation renders each candidate independently using the single-candidate synthesizer.
On the same M4 Mini hardware this achieves ${\sim}200$ evaluations per second---$2.8\times$ slower than the batched approach due to inter-process communication overhead and the inability to share FFT plans across workers.

\paragraph{Cloud comparison.}
We also benchmark on an Azure Standard\_D48s\_v5 VM (48 vCPU, Intel Xeon Platinum 8370C).
With 48-way multiprocessing the cloud VM achieves ${\sim}220$ evaluations per second.
The M4 Mini's unified memory architecture and high single-core performance deliver ${\approx}2.5\times$ higher per-core throughput than the cloud VM, making consumer Apple Silicon competitive with datacenter hardware for this workload.

A complete 100K-evaluation CMA-ES run takes approximately 3 minutes on the M4 Mini with the batched synthesizer.
